% 蚂蚁技术研究院 密码学实验室 二面（老板面）—— 自用小抄，不共享屏幕。
% Build: xelatex ant-round2-crib.tex (twice).
% 事实来源：SHIELD = PaperSrc/682277aee485710d428ea960/SP/Bodies/
%          masking = PaperSrc/6ab6c47fcc98e0b1b3b2d1c6/src/bodies/
%          HERO    = PaperSrc/6984fb4d3ed4867cd184e11d/NDSS_version/
%          简历    = resume/Changqi_Sun_AntGroup_PQC2027.pdf（v2-crypto，已提交，不改）
%          文献    = Celi et al. ePrint 2026/013；Quorus ePrint 2025/1163；Trilithium 2025/675；Tang et al. 2024/1067
% 标【待确认】的是推断或没核实的，面试前自己过一遍。
\documentclass[aspectratio=169,10pt,fontset=none]{ctexbeamer}
\setCJKmainfont{PingFang SC}[BoldFont={PingFang SC Semibold}]
\setCJKsansfont{PingFang SC}[BoldFont={PingFang SC Semibold}]
\usetheme{default}
\setbeamertemplate{navigation symbols}{}
\setbeamertemplate{footline}[frame number]
\usepackage{booktabs,tabularx,array,stmaryrd}
\usepackage{tikz}
\usetikzlibrary{positioning,arrows.meta,fit}
\definecolor{accent}{HTML}{1F4E79}
\definecolor{warn}{HTML}{B22222}
\definecolor{ok}{HTML}{1F6B3A}
\setbeamercolor{frametitle}{fg=accent}
\setbeamercolor{structure}{fg=accent}
\setbeamerfont{frametitle}{size=\large,series=\bfseries}
\newcommand{\warn}[1]{{\color{warn}\textbf{【待确认】}#1}}
\newcommand{\say}[1]{{\color{ok}\textbf{口径：}}#1}
\newcommand{\src}[1]{{\scriptsize\color{gray}#1}}
\newcolumntype{Y}{>{\raggedright\arraybackslash}X}
\setlength{\parskip}{2pt}

\title{蚂蚁 · 密码学实验室 二面小抄}
\subtitle{老板熟 MPC：SHIELD 的「需求 $\to$ 选择」故事 + masking + ML-KEM}
\author{自用，不共享屏幕}
\date{投递简历：v2-crypto / Changqi\_Sun\_AntGroup\_PQC2027.pdf（不改，口头更正）\quad AI Coding 笔试：暂不做}

\begin{document}
\footnotesize

\begin{frame}
\titlepage
\vspace{-1em}
\begin{itemize}
  \item 页序：一面回顾 $\to$ 开场 $\to$ \alert{纠错总表} $\to$ \textbf{1 SHIELD} $\to$ \textbf{2 Falcon/FLOCK} $\to$ \textbf{3 masking} $\to$ \textbf{4 HERO} $\to$ \textbf{5 threshold ML-KEM} $\to$ \textbf{6 概念讲解 E1--E6} $\to$ 术语 + 数字速查
  \item {\color{warn}红字【待确认】}= 推断或没核实，面试前自己扫一遍
\end{itemize}
\end{frame}
\begin{frame}{一面回顾（王宇辰，11-04，近 2 小时）}
\begin{columns}[T]
\column{.5\textwidth}
\textbf{被问到的}
\begin{itemize}
  \item HERO 讲了很多（按时间顺序讲的，没问题）
  \item SHIELD：\textbf{为什么 malicious？为什么 SPDZ？为什么 WRK17？}\\
        \quad 当时答「导师推荐 / SPDZ 好做 malicious」—— \alert{不够 convincing}
  \item \textbf{每个人都问性能}：最终能做到多少 latency、几方
  \item 面试官提到 \textbf{ISW} —— 我没接住（见 E1）
  \item 简历第 1 条 masking + 省内存：还没投，要能说清状态
\end{itemize}
\column{.5\textwidth}
\textbf{对方透露的}
\begin{itemize}
  \item 老板熟 MPC，会追问设计取舍
  \item 组里想做 \textbf{threshold ML-KEM}，但自己也不确定能不能做 $\Rightarrow$ 我的定位：\textbf{可行性探索}（Part 5）
  \item 他提到 threshold ML-DSA 有好几篇新工作 —— 见「相关工作」页
\end{itemize}
\textbf{本轮要纠正的记忆}
\begin{itemize}
  \item 最多测到 \textbf{5 方}（不是 6/8；6 是 Celi/Borin 的上限）
  \item 「< 500 ms」是\textbf{单次迭代}；87@10 ms RTT 是 530 ms
\end{itemize}
\end{columns}
\end{frame}
\begin{frame}{开场 30 秒（中文说，术语英文）}
\begin{itemize}
  \item VT CS 四年级 PhD，导师 Thang Hoang，方向 applied crypto，2028-06 毕业。数学本科。
  \item 主线：\textbf{lattice 签名的 threshold 化}，而且坚持\alert{输出就是标准签名}：
  \begin{itemize}
    \item SHIELD（S\&P 2026 共同一作 + 报告人）：threshold ML-DSA，MPC 里跑\textbf{不改动}的 FIPS 204
    \item FLOCK（一作，USENIX Sec 2027 在投）：Falcon-compatible threshold signature
  \end{itemize}
  \item 实现侧：masked ML-DSA on Cortex-M4 —— 发现寄存器级防护在 bus / port 上漏，并把峰值 RAM 降 37--42\%
  \item 另有 ORAM 工作 HERO（一作，NDSS 2027 在审）
  \item 一句收尾：\textbf{「我做的事情都是：在不改标准算法的前提下，把它放进更苛刻的环境里（多方 / 侧信道 / 小内存）」} —— 正好接 threshold ML-KEM
\end{itemize}
\end{frame}
\begin{frame}{先读这页：我之前说错 / 记错 / 没意识到的（逐条改口）}
{\small\renewcommand{\arraystretch}{1.2}
\begin{tabularx}{\textwidth}{@{}p{4.4cm}Y@{}}\toprule
\textbf{我说过 / 记得的} & \textbf{正确说法（怎么说出口）} \\\midrule
「masking 本质上就是 MPC in the head」 & \alert{MPC-in-the-head 是用 MPC 造 ZK proof}（IKOS07）。masking 应说「\textbf{单设备上的 MPC}：ISW private circuits，probing model 下每个中间值拆成 $d+1$ 份」—— 见讲解页 E1/E2 \\
「WLK30 / WRK 那篇」 & \textbf{WRK17} = Wang--Ranellucci--Katz, ``Global-Scale Secure Multiparty Computation'', CCS 2017（authenticated garbling 的多方版）\\
「最多做到 6 或 8 方」 & SHIELD 测了 \textbf{2--5 方}。6 方是 Celi et al. 的上限 \\
「online < 500 ms」 & \textbf{单次成功迭代} 198--536 ms（87@10 ms RTT 超 500）；端到端还要乘期望重试 $\approx$4 次 \\
「Efficient threshold ML-DSA < 100 ms」 & Celi et al.（USENIX Sec'26）本机 0.5--12 ms、LAN < 20 ms、\textbf{跨洲 WAN 0.6--0.75 s}；「sub-100 ms」出自未评审的 Kao arXiv \\
「double shares，每方存两份」 & \textbf{short replicated secret sharing}；每方存 2 份只是 $(N,T)=(3,2)$ 的特例，一般是 $\binom{N-1}{N-T}$ 份 \\
「用 SPDZ 是因为比较容易实现 malicious」 & SPDZ 用 \textbf{IT-MAC} 在 \textbf{dishonest majority} 下给 active security with abort —— 关键词是 dishonest majority，不是「容易」\\
\bottomrule
\end{tabularx}}
\end{frame}

\begin{frame}{先读这页（续）：没意识到的限制 + 术语}
{\small\renewcommand{\arraystretch}{1.2}
\begin{tabularx}{\textwidth}{@{}p{4.4cm}Y@{}}\toprule
\textbf{我说过 / 记得的} & \textbf{正确说法（怎么说出口）} \\\midrule
「用它是因为导师推荐」 & 按 R1--R4「需求 $\to$ 选择」讲（SHIELD 第 2 页）\\
「SHIELD 抗 $n-1$ 恶意」（没意识到的限制）& \alert{online 才是}；offline 的 triples / daBits 由\textbf{可信 dealer} 生成 \\
简历「890K$\to$355K vs SCALE-MAMBA」 & 口头说「相对我们最初的朴素编码」；和 SCALE-MAMBA 比的是加法器 375$\to$294 \\
masking：「MakeHint / Keccak 都是 ISW 掩码」 & MakeHint \textbf{明文算}（检查全过并合并后）；Keccak 用 \textbf{TI}（threshold implementation）不是 ISW；论文\textbf{不声称} SNI/PINI \\
「TI 的 threshold」 & \alert{和 threshold signature 的 threshold 无关}：TI 指「每个输出份只依赖部分输入份」（non-completeness）\\
「strict ML-KEM」 & 是 \textbf{threshold ML-KEM}；组里也不确定能不能做 $\Rightarrow$ 我的定位是「可行性探索」\\
「Efficient Threshold ML-DSA = ePrint 2025/1166」 & 1166 已撤回、拆分；现在是 \textbf{ePrint 2026/013}（另一半 Threshold Raccoon 是 2026/419）\\
\bottomrule
\end{tabularx}}
\end{frame}
\begin{frame}
\vfill\centering{\Huge\color{accent} Part 1 · SHIELD}\\[6pt]{\normalsize threshold ML-DSA（S\&P 2026）—— 老板重点}\vfill
\end{frame}
\begin{frame}{SHIELD 是什么（一页讲清）}
\centering
\begin{tikzpicture}[font=\scriptsize, >=Stealth,
  box/.style={draw, rounded corners, align=center, inner sep=3pt, text width=3.1cm, minimum height=1.0cm}]
\node[box, fill=gray!10] (a) {可信 CA 一次性 setup\\ 生成 ML-DSA 密钥，\\ 分发 SPDZ/WRK17 shares};
\node[box, fill=accent!10, right=.5cm of a, text width=4.6cm] (b) {$n$ 个 CA 在 MPC 里联合跑\\ \textbf{未修改的} FIPS 204 Sign\\ （含拒绝采样、SHAKE、MakeHint）};
\node[box, fill=ok!12, right=.5cm of b] (c) {输出 = 标准 ML-DSA 签名\\ 标准 verifier 直接验\\ \textbf{bit-exact}};
\draw[->] (a) -- (b); \draw[->] (b) -- (c);
\end{tikzpicture}
\vspace{.4em}
\begin{itemize}
  \item 场景：\textbf{FPKI}（联邦 PKI）—— 多个 CA 共同签发证书，任何单个 CA 都拿不到可用私钥
  \item 威胁模型：malicious、static、QPT 敌手，最多腐化 $n-1$ 个 CA；security-with-abort；quantum-UC 证明
        \src{model.tex:47, security\_analysis.tex:52}
  \item 门限：\textbf{$n$-of-$n$}（$t=n$）—— 「所有 CA 都参与，最大化 trust 和 accountability」\src{proposed\_method.tex:2}
  \item 引擎：\textbf{SPDZ}（算术：NTT、$A\mathbf y$、$c\mathbf s_1$）+ \textbf{WRK17}（布尔：SHAKE/Keccak-f、ExpandMask、Decompose/HighBits、SampleInBall、CheckNorm、MakeHint）+ \textbf{daBits} 跨域桥接
  \item 实现：SPDZ 自己从零写；WRK17 用 EMP-toolkit；布尔电路用 Verilog $\to$ Yosys $\to$ Bristol format
  \item 定位一句话：前人做到 \emph{functional interchangeability}（标准 verifier 能验），我们是 \alert{algorithmic FIPS-204 compliance}（算法本身没改）\src{background.tex:17}
\end{itemize}
\end{frame}
\begin{frame}{为什么这么选：先讲需求，再推出选择（核心页）}
\renewcommand{\arraystretch}{1.25}
\begin{tabularx}{\textwidth}{@{}p{3.6cm}Y p{3.3cm}@{}}
\toprule
\textbf{需求 / 约束} & \textbf{推导} & \textbf{选择} \\ \midrule
\textbf{R1} 持钥方本身就是敌手（CA 各自独立、互不信任） &
半诚实不够：恶意方在 $\mathbf z$ 的份额或 reject bit 上加 $\delta$，就能让本该被拒的 $\mathbf z$ 被打开 —— 而\textbf{拒绝采样正是 ML-DSA 藏 $\mathbf s_1$ 的机制}；还有 selective failure、nonce 偏置 &
\textbf{malicious security}（IT-MAC 认证每个中间值）\\
\textbf{R2} $n$-of-$n$，要容忍 $n-1$ 个坏 CA &
这就是 dishonest majority $\Rightarrow$ 排除 Shamir / honest-majority RSS（Quorus 要 $n=2t+1$，3 方只容 1 个坏人）&
\textbf{dishonest-majority MPC} \\
\textbf{R3} 输出必须和 FIPS 204 \textbf{逐比特一致}，不改算法 &
排除「每方本地拒绝采样」(Celi：换成 hyperball 分布)、「加噪声 $e_w$」(Quorus)、「hash 放 MPC 外 + 同态承诺」(Tang) $\Rightarrow$ 只能把\textbf{整个 Sign 原样搬进通用 MPC} &
\textbf{通用 MPC}，不是定制协议 \\
\textbf{R4} ML-DSA 是混合计算：$\mathbb Z_q$ 多项式算术 + 大量位运算 &
全 SPDZ：比较 / 位分解在算术域要 $O(\log q)$ 轮；全 GC：23-bit 模乘上千 AND 门 $\Rightarrow$ 各做擅长的 &
\textbf{SPDZ + WRK17 + daBits} \\
\bottomrule
\end{tabularx}
\vspace{.4em}

\say{「最初是导师建议的组合，但我回头复盘，它其实是这四条约束推出来的唯一自然解；如果放宽 R2 或 R3，就会走向 Quorus 或 Celi 那样更快的路线。」}
\end{frame}
\begin{frame}{为什么 SPDZ、为什么 WRK17（被追问细节时）}
\begin{columns}[T]
\column{.5\textwidth}
\textbf{SPDZ}（Damgård et al. 2012）
\begin{itemize}
  \item 份额 $\llbracket x\rrbracket_i=(x_i,\ m_i)$，$\sum m_i=\alpha x$，$\alpha$ 是全局 MAC key
  \item 线性运算本地免费；乘法用 Beaver triple，只需打开 $d,e$
  \item online 信息论安全，开销约半诚实版的 2 倍
  \item 任意 $N$，容忍 $N-1$ 恶意 —— 正好满足 R2
  \item ML-DSA 的 $q=8380417$ 是素数 $\Rightarrow$ 主体就是 $\mathbb F_q$ 线性代数，和 SPDZ 天然契合
  \item 联合拒绝（不是每方各自拒绝）$\Rightarrow$ 期望重试次数 = 单机 ML-DSA（4.25 / 5.1 / 3.85）
\end{itemize}
\column{.5\textwidth}
\textbf{WRK17}（Wang–Ranellucci–Katz, CCS'17, ``Global-Scale MPC''）
\begin{itemize}
  \item authenticated garbling：预处理生成带 IT-MAC 的 AND triple（TinyOT 风格）
  \item malicious、dishonest majority（$N-1$），\textbf{online 常数轮}（我们是 4 轮），与电路深度无关 \src{experiment.tex:64}
  \item 用在：SHAKE/Keccak-f、ExpandMask、Decompose/HighBits、SampleInBall (Fisher–Yates)、CheckNorm、MakeHint、A2B 加法器 \src{design\_rational.tex:15--64}
  \item 电路优化：SampleInBall 890K $\to$ 355K gates
        \warn{简历写成「vs SCALE-MAMBA」；论文里是相对\textbf{我们最初的朴素编码}。口头说「相对最初编码」。和 SCALE-MAMBA/EMP 比的是 32-bit 加法器 375 $\to$ 294 门（$-27\%$）}
\end{itemize}
\textbf{daBits}：同一个随机比特在算术域和布尔域各有一份认证 share（Rotaru–Wood, Marbled circuits）$\Rightarrow$ A2B / B2A 只需打开一次掩码值
\end{columns}
\end{frame}
\begin{frame}{性能：数字要说准}
\begin{columns}[T]
\column{.55\textwidth}
\textbf{单次成功迭代的 online latency（ms），$n=2\to5$}\\
{\scriptsize 本地集群 + 模拟 RTT：Low$\approx$1 ms / Med$\approx$3 ms / High$\approx$10 ms，10 Gbps}\\[2pt]
\begin{tabular}{@{}lccc@{}}\toprule
 & Low & Med & High \\\midrule
ML-DSA-44 & 198--205 & 253--258 & 414--418 \\
ML-DSA-65 & 246--252 & 294--298 & 455--461 \\
ML-DSA-87 & 321--326 & 369--376 & \alert{530--536} \\\bottomrule
\end{tabular}
\src{experiment.tex:263--276}
\vspace{.4em}

\textbf{每个 CA 的 online 开销}：23 轮；7.7 / 10.9 / 14.7 MB（+ 少量 $\propto n$）；内存 160--280 MB\\
\textbf{吞吐}：16 核并行，44@Low $\approx$20 sig/s，44@High $\approx$6 sig/s
\column{.45\textwidth}
\textbf{怎么说（主动讲清口径）}
\begin{itemize}
  \item 「单次迭代 online 0.2--0.33 s（LAN），10 ms RTT 下 ML-DSA-44/65 不到 0.46 s；\textbf{最多测到 5 方}」
  \item 「端到端要乘期望重试 $\approx$4 次；可以像 Celi 那样并行跑 $K$ 份来换 latency」
  \item 「没测真实 WAN；23 轮 online 在 100 ms RTT 下光等待就 2.3 s/迭代 —— 这是 SPDZ 路线的代价」
  \item \textbf{offline 很重}：WRK17 预处理是瓶颈；摊销后一次成功迭代按\alert{分钟}计（44@Low 2.8--8.6 min）\src{perf\_full\_table.tex:139--153}
\end{itemize}
\textbf{对比}：Celi 本机 0.5--12 ms、LAN < 20 ms、跨洲 WAN 0.6--0.75 s；Quorus online 计算 6--343 ms（不含网络）
\end{columns}
\end{frame}
\begin{frame}{软肋 Q\&A（老板最可能追问，主动承认 + 给改进）}
\renewcommand{\arraystretch}{1.2}
\begin{tabularx}{\textwidth}{@{}p{3.4cm}Y@{}}
\toprule
\textbf{追问} & \textbf{回答} \\\midrule
offline 谁做？ & SPDZ 的 triples / daBits / AGV 由\textbf{可信 dealer}（setup 那个 CA）生成，所以 dishonest majority \textbf{只覆盖 online}。改进：MASCOT / Overdrive(LowGear) 生成 triples，edaBits（Escudero et al. CRYPTO'20）；WRK17 预处理本来就是各方自己跑的 \\
MAC soundness？ & $\mathbb F_p$ 上单次 check 伪造概率 $\approx 1/p$。若 $p=q\approx2^{23}$ 只有 $2^{-23}$ \warn{论文只写 $\mathbb F_p$；A2B 按 $|q|=23$ bit 计，推测 $p=q$。回去查代码确认}。修法：$k$ 个独立 MAC key（$2^{-23k}$）或扩域 MAC \\
WAN / 轮数？ & 23 轮 online，RTT 敏感。Celi 每次尝试 3 轮；这是「不改算法」的代价 \\
后量子？ & quantum-UC 只覆盖 IT-MAC 部分；WRK17 / EMP 的 base OT 通常是 ECC，不是 PQ。修法：base OT 换 PQ OT，OT extension 本身是对称密码 \\
abort 能识别谁吗？ & 不能。只有 security-with-abort，static corruption。identifiable abort 要额外机制（SPDZ-IA 类） \\
被拒迭代的 opening？ & \warn{自查：被拒迭代里打开过的值是否也进 batch MAC check？selective failure 要求「影响后续行为的值」在用之前被认证} \\
并发安全（加分）& 已处理 Kyster et al.（S\&P'25 ``Rushing at SPDZ''）：并发下 MAC key 可被提取 $\Rightarrow$ 对所有 partial openings 做一次 batch MAC check \src{design\_rational.tex:75--77} \\
\bottomrule
\end{tabularx}
\end{frame}
\begin{frame}{相关工作：threshold ML-DSA 地图}
{\scriptsize
\renewcommand{\arraystretch}{1.15}
\begin{tabularx}{\textwidth}{@{}p{2.9cm}p{1.6cm}Y p{2.3cm}p{1.1cm}Y@{}}
\toprule
\textbf{工作} & \textbf{出处} & \textbf{技术} & \textbf{安全} & \textbf{方数} & \textbf{性能} \\\midrule
\textbf{Celi et al.}「Efficient Threshold ML-DSA」(= NIST MPTC ``Mithril'') & USENIX Sec'26, ePrint 2026/013 & \textbf{short RSS}，每方本地 hyperball 拒绝；每次尝试 3 轮 & malicious, dishonest maj.; game-based; abort 不可识别 & $\le 6$（8 勉强）& 本机 0.5--12 ms；LAN < 20 ms；WAN 27 ms（亚洲内）/ 0.6--0.75 s（跨洲）\\
\textbf{Quorus} (Bienstock et al., JPMC) & USENIX Sec'26, 2025/1163 & Shamir mod $q$ + edaBits；加噪声 $e_w$ 使公开被拒 $w_1$ 安全 & malicious, \textbf{honest maj.}; UC & 任意（测到 63）& online 计算 6--343 ms，不含网络；16--29 轮/次 \\
\textbf{Trilithium} & ESORICS'26, 2025/675 & 2PC + BeDOZa MAC + 诚实的关联随机数提供方 & 需 CRP 诚实 & 2 & 秒级；CRP 每次 55 MiB \\
\textbf{Tang et al.} \alert{（最接近）} & TIFS'23, 2024/1067 & \textbf{$t$-of-$n$ SPDZ} + 分布式拒绝采样；hash 在 MPC 外 + 同态承诺 & 主动安全 & 测到 12 & 0.5 s (2 方)，7.3 s (12 方) \\
\textbf{SHIELD（我们）} & S\&P'26 & SPDZ + WRK17 + daBits，整个 Sign 原样进 MPC & malicious, dishonest maj.; quantum-UC & 2--5 & 单次迭代 0.2--0.54 s \\
\bottomrule
\end{tabularx}}
\vspace{.3em}

\say{「和 Tang 的区别：他们 hash 放 MPC 外、用同态承诺，输出不是精确 ML-DSA；我们 SHAKE 也在 MPC 里，bit-exact。和 Celi 比我们慢两个数量级，但他们改了拒绝分布、只到 6 方。」}\\
\src{不要拿 Kao (arXiv 2601.20917) 做对比 —— 未评审，同作者的 TALUS 已被 Niot (ePrint 2026/1386) 攻破}
\end{frame}
\begin{frame}{Celi et al. 的「double shares」= short replicated secret sharing}
\begin{columns}[T]
\column{.52\textwidth}
\textbf{怎么分}
\begin{itemize}
  \item 对每个大小 $N-T+1$ 的子集 $I$，采一个\textbf{短的} $\mathbf s_I\sim\chi_s$；$\mathbf s=\sum_I \mathbf s_I$
  \item 参与者 $i$ 持有所有 $i\in I$ 的 $\mathbf s_I$
  \item 总份数 $\binom{N}{T-1}$；每方 $\binom{N-1}{N-T}$ 份
  \item \textbf{(3,2)：共 3 份，每方 2 份} ← 「每人存两份」；(6,4)：共 20 份，每方 10 份
  \item 至多 $T-1$ 方被腐化 $\Rightarrow$ 至少一个 $\mathbf s_I$ 全在诚实方手里
\end{itemize}
\textbf{为什么必须「短」}：Shamir 份额在 $\mathbb Z_q$ 上均匀、Lagrange 系数大 $\Rightarrow$ 部分签名范数大，没法本地拒绝。RSS 重构系数是 0/1 $\Rightarrow$ 每方部分秘密本身就短
\column{.48\textwidth}
\textbf{怎么签}
\begin{itemize}
  \item 每方本地跑 FS-with-aborts：$\mathbf z_i=c\,\mathbf s_i^{\mathrm{part}}+\mathbf r_i$，加起来就是签名
  \item 各方都要通过 $\Rightarrow$ 成功率 $\approx p^T$ $\Rightarrow$ 换成 \textbf{hyperball 拒绝}（$\ell_2$，像高斯一样紧），imbalanced 半径（hint 只作用在 $\mathbf z^{(2)}$）
  \item 两阶段：各方本地拒绝（安全必需）+ combine 时全局检查（尺寸兼容）
  \item 承诺用完整 MLWE 样本 $\mathbf w_i=A\mathbf r_i$，被拒也能安全公开
  \item 3 轮/次：commit $H(\mathbf w_i)$（防 rushing）$\to$ 公开 $\mathbf w_i$ $\to$ 发 $\mathbf z_i$；$K=\lceil1/p\rceil$ 份并行
\end{itemize}
\textbf{和我们的本质差别}：输出能过标准 verifier，但\alert{签名分布不是 FIPS 204 的}；份数随 $N$ 指数涨 $\Rightarrow$ $N\le6$
\end{columns}
\end{frame}
\begin{frame}{如果现在重做 SHIELD（加分页）}
\begin{itemize}
  \item \textbf{offline 去掉可信 dealer}：triples 用 MASCOT/LowGear，daBits $\to$ edaBits（一次拿一整串比特，A2B 更省）
  \item \textbf{latency}：并行跑 $K$ 个候选（Celi 的 $K=\lceil1/p\rceil$ 思路），一次成功率 $\ge 1-e^{-1}$；批量 MAC check 合并所有候选
  \item \textbf{$t$-of-$n$}：SPDZ 本身支持 $t$-out-of-$n$ 变体（Tang 那条路线）；FPKI 之外的托管 / 钱包场景需要它
  \item \textbf{MAC soundness}：多 MAC key 或扩域，让统计安全到 40--64 bit
  \item \textbf{场景拆分}：如果能接受 honest majority（例如同一机构的 3 台 HSM），直接走 Quorus / Shamir 更快；如果 $N\le6$ 且不要求分布一致，走 Celi
  \item 一句话：\textbf{「SPDZ+WRK17 是在最强约束下的解；约束一放松，就该换更快的定制协议。」}
\end{itemize}
\end{frame}
\begin{frame}
\vfill\centering{\Huge\color{accent} Part 2 · Falcon}\\[6pt]{\normalsize threshold-Falcon artifact 修复 + FLOCK}\vfill
\end{frame}
\begin{frame}{Falcon 两条（简历第 2 条 + SUMMARY 里的 FLOCK）}
\begin{columns}[T]
\column{.5\textwidth}
\textbf{threshold-Falcon artifact 修复（merged PR）}
\begin{itemize}
  \item 复现 CCS'26 ``Thresholdizing Standardized FALCON'' 时，签名\textbf{连续}多次 l2bound 失败 —— 按 Falcon 的拒绝率统计上不可能 $\Rightarrow$ 定位到 sampler
  \item 根因：emp-tool \texttt{PRG::random\_data} 字节/元素单位不匹配 $\Rightarrow$ sampler 缓冲区 \textbf{7/8 读的是未初始化堆内存}，签名范数偏大 $\approx$15\%
  \item 第二个 bug：\texttt{read()/write()} 返回值存进 \texttt{size\_t}，$-1$ 变 SIZE\_MAX，I/O 流静默失步
  \item 修后：均方范数与期望差 0.4\%，零拒绝，受影响的 run 快 $\sim$9$\times$；作者合并
\end{itemize}
\column{.5\textwidth}
\textbf{FLOCK（一作，USENIX Sec'27 在投）}
\begin{itemize}
  \item 难点：Falcon 的 ffSampling 是浮点、$2n$ 次顺序采样、每次以秘密为中心 $\Rightarrow$ 前人 threshold Falcon 每签名上万轮
  \item (1) 换 sampler：Peikert offline/online + Zalcon + Antrag trapdoor $\Rightarrow$ 纯整数、可并行，顺序深度 2
  \item (2) 「以秘密为中心」的采样放到 MPC 外：打开 masked center $c+e$，各方明文采样，回到 MPC 只做线性运算
  \item 结果：9 轮 online；4 方 LAN 25 ms vs Garg et al. 24.2 s；签名过 Falcon verifier（$-4/12$ forgery bits，$+9/32$ bytes）
  \item 对外说 \textbf{Falcon-compatible}，不说 ``standard Falcon''
\end{itemize}
\end{columns}
\end{frame}
\begin{frame}
\vfill\centering{\Huge\color{accent} Part 3 · masking}\\[6pt]{\normalsize Clean Registers, Leaky Buses（在写，S\&P 2027）}\vfill
\end{frame}
\begin{frame}{masking 总览：同一个 ML-DSA，从「多台机器」换成「一块芯片」}
\begin{columns}[T]
\column{.5\textwidth}
\textbf{high-level 一句话}：「SHIELD 把 ML-DSA 拆给多台机器算；masking 是\textbf{在一块 Cortex-M4 上把每个中间值拆成几份}，probing model 下也是一种 MPC（ISW private circuits）。两边的难点一模一样：\alert{算术域和布尔域之间的转换}。」\\[3pt]
\textbf{论文}：``Clean Registers, Leaky Buses: Hardening Masked ML-DSA under Memory Constraints''，基于 mldsa-native 的一阶掩码签名路径；目标 S\&P 2027 Cycle 2（摘要 11-10 / 全文 11-17）\\[3pt]
\textbf{只合并（unmask）这几样}：每次尝试的 $w_1$、$\tilde c$；每行范数检查 1 个判定位；$\mathbf r,\mathbf z$ 只在三项检查都通过后才合并；\textbf{被拒尝试的 $\mathbf z,\mathbf r$ 永不合并}（FS-with-aborts 只保证被接受的那组可模拟）
\column{.5\textwidth}
\textbf{份额数（每个都有理由）}\\[2pt]
{\scriptsize\renewcommand{\arraystretch}{1.15}
\begin{tabularx}{\linewidth}{@{}p{2.3cm}cY@{}}\toprule
域 & 份数 & 理由 \\\midrule
算术 mod $q$（$\mathbf y,\mathbf w,\mathbf z,\mathbf r,\hat{\mathbf s}_{1,2}$）& 2 & 一阶；$A\mathbf y$、$c\mathbf s_1$ 对秘密线性，$A,c$ 公开 $\Rightarrow$ 逐份额算，不要随机数 \\
布尔：Keccak + $\mathbf y$ 采样器 & 4 & 针对 Weakest Link（Hermelink et al.）打 $\mathbf y$；他们建议升到 2 阶，\textbf{4 份是我们自己的选择} \\
布尔：Decompose + 两个范数检查 & 3 & Balasch et al.：$n$ 份在 transition 泄漏下只有 $\lfloor(n-1)/2\rfloor$ 阶 $\Rightarrow$ 3 是一阶的最小值；代价 +97.46\% 指令 \\
$t_0$ & 0 & FIPS 204 §6.1：不必当秘密 \\\bottomrule
\end{tabularx}}
\end{columns}
\end{frame}
\begin{frame}{A2B / B2A：两个域怎么转（你关心的那块）}
\begin{columns}[T]
\column{.5\textwidth}
\textbf{B2A：ExpandMask $\to\mathbf y$}（4 布尔份 $\to$ 2 算术份 mod $q$）
\begin{itemize}
  \item 做法（BBEFG18 §5「减随机份额再合并」）：抽随机算术份 $s$ $\to$ 在 4 个布尔份上用 SecAdd 算 $a_0=\gamma_1-z-s \bmod q$（再一次 csubq）$\to$ \textbf{只合并 $a_0$}，$(a_0,s)$ 就是算术份
  \item SecAnd：4 份非完备 TI，3 个 fresh 字
  \item 备选（代码里非默认）：SPOG19 按位 B2A$_q$，每系数 54 次乘法
\end{itemize}
\textbf{A2B$_q$：Decompose / HighBits}（2 算术份 $\to$ 3 布尔份）
\begin{itemize}
  \item Barthe et al. 18：每个算术份各自布尔拆分，用 \textbf{Kogge--Stone SecAdd} 相加，再 $n-1$ 次 csubq（mod $q$ 不是 $2^k$，所以要条件减）
  \item HighBits 在布尔份上算（ISW AND + Kogge--Stone；乘常数用移位加法链）；只合并 3 个布尔份得 $w_1$；$w_0=w-2\gamma_2 w_1$ 留在算术域
\end{itemize}
\column{.5\textwidth}
\textbf{范数检查 MChkNorm$_3$}：份 0 加 $B-1$ $\to$ A2B $\to$ 加公开常数取最高借位 $\to$ ISW AND 累积各行 $\to$ \textbf{只 XOR 合并最后 1 bit}\\[3pt]
\textbf{没掩码 / 没用的}（别说错）
\begin{itemize}
  \item \alert{MakeHint 明文算}（检查全过并合并后）；没有 1-bit B2A
  \item Goubin A2B 没用：值模型成立，\textbf{transition 模型不成立}
  \item Keccak 的 $\chi$：4 份 TI + changing of the guards（不是 ISW/DOM）；随机数 28.8 KB $\to$ 2.9 KB / 置换
  \item 不声称 SNI / PINI，只有一阶、值模型的逐 gadget 论证
\end{itemize}
\textbf{代价}：Kogge--Stone 每次加法 10 SecAnd + 9 refresh（ripple 要 31+31）；真板每次尝试 $\approx$27k 次 Keccak-f、3.7 MB 随机数，掩码 RNG 占 65.9\% 时间\\[3pt]
\textbf{和 SHIELD 对照}：SHIELD 用 daBits 一次打开掩码值做转换（MPC 里可以「打开」）；芯片上没人能安全「打开」，只能用 SecAdd 电路在份额上算 $\Rightarrow$ 贵
\end{columns}
\end{frame}
\begin{frame}{泄漏发现：寄存器干净，总线漏}
\begin{columns}[T]
\column{.52\textwidth}
\textbf{机制}：2 份 refresh 写成 \texttt{a0\^{}=r; a1\^{}=r}，两份从没先后进同一个寄存器，但\textbf{先后经过数据总线}；布尔份 $\mathrm{HD}(x_0,x_1)=\mathrm{HW}(x)$ $\Rightarrow$ 直接漏明文的 HW
\begin{itemize}
  \item 同一个二进制：寄存器模型 $\max|t|\le5.83$（干净）；总线模型 chknorm / decompose $|t|=199.00$
  \item \textbf{R1}：同一秘密的两份不能在同一个有状态位置（寄存器 / 总线 / 操作数端口）相邻
  \item \textbf{R2}：每个 barrier 必须清并声明\textbf{所有}可分配寄存器
\end{itemize}
\textbf{-O2 vs -O3}：barrier 只清 caller-saved r0--r3/r12，-O2 下编译器把相邻份放进 callee-saved r6/r8 $\Rightarrow$ decompose 24 点超阈，$\max|t|$ 14.10；改全寄存器 barrier（r0--r12）后消失
\column{.48\textwidth}
\textbf{对策}：bus-and-port barrier（栈字写 0 再读回 + 全 0 寄存器跑 \texttt{mla} 驱动端口）+ 3 份检查路径 + 全寄存器 barrier；+14.65\% 指令\\[3pt]
\textbf{评估}：Cortex-M4 仿真，register / bus / operand 三种模型的一、二阶 TVLA，全签名（最多 2.675 亿条指令 / trace），CPA + template，各带 positive control；TVLA 阈值随 trace 长度：$\sqrt{2\ln(2m/\alpha)}$，$\alpha=0.05$\\[3pt]
\textbf{诚实说的限制}
\begin{itemize}
  \item \alert{没上真机测泄漏}；对两个公开 M4 硬件数据集检出 27/46、26/37
  \item 3 个 FAIL 原样记录：template GE 15.97 < 16；种子 Keccak 1 点 6.14 > 6.08（最多漏份 1）；-O2 rhoprime 序言/尾声
  \item 只做一阶；没做水平攻击（LAZY\_Y 让同一 $\mathbf y$ 算 3 次）
\end{itemize}
\end{columns}
\end{frame}
\begin{frame}{省内存的逻辑（ML-DSA-44：50,612 $\to$ 31,828 B，$-37.1\%$）}
\begin{columns}[T]
\column{.6\textwidth}
{\scriptsize\renewcommand{\arraystretch}{1.12}
\begin{tabularx}{\linewidth}{@{}p{1.9cm}Y r r@{}}\toprule
步骤 & 做法 / 为什么安全 & $\Delta$ & 峰值 \\\midrule
d 默认掩码 & — & & 71,900 \\
n0 加固+LAZY & LAZY\_A：$\hat A$ 逐元素现算（$\rho$ 公开，$-$15,264 B）；LAZY\_Y：$\mathbf y$ 从 $\langle\rho''\rangle_4$ 重采（$-$6,072 B，\alert{但同一 $\mathbf y$ 多次暴露}）& & 50,612 \\
M1 & 公开的 $t_0$ 打包存 & $-$2,424 & 48,188 \\
M6a & 删掉掩码路径从不读的 NTT($\mathbf y$) 暂存区 & $-$3,072 & 45,116 \\
M6b & 明文 $\mathbf r$ 写进已死的掩码 $\mathbf w$ 存储 \warn{覆写曾存份额的字：字节粒度写 0 在一个设备上漏，未排除} & $-$4,096 & 41,020 \\
M7 & 明文 $\mathbf z$ 行放进 $\hat A$ 槽位（生命周期不重叠，公开盖公开）& $-$1,024 & 39,996 \\
M3 & $\hat{\mathbf s}_{1,2}$ 的份 1 换成 32 B 种子（SHAKE128 展开）& $-$8,168 & \textbf{31,828} \\\bottomrule
\end{tabularx}}
\src{65：66,244 $\to$ 40,100（$-39.5\%$）；87：84,948 $\to$ 49,396（$-41.9\%$）。n0$\to$n3 只多 +0.21\% 指令}
\column{.4\textwidth}
\textbf{思路三类}（面试就讲这三句）
\enumerate
  \item \textbf{能重算就不存}：公开的 $\hat A$ 现算；$\mathbf y$ 重采（拿暴露换内存）
  \item \textbf{按生命周期复用}：死缓冲删掉，生命周期不重叠的缓冲共用
  \item \textbf{份额压成种子}：一份份额 = PRF(seed)，密钥对象 $-$8 KB
\endenumerate
\textbf{M3 的细节}：每系数取 48 位定长 Montgomery 约简（拒绝采样会让 trace 长度随密钥变 $\Rightarrow$ 弃用）；偏差 $2^{-30.45}$；refresh 先减新份再加旧份（反过来份 0 会短暂等于明文）；新假设 SHAKE128 是 PRF\\[2pt]
\textbf{没做 / 前人}：份额打包到 23 位（再省 $\approx$2.3 KB）、按行重算 $\mathbf w$（省 8 KB 但每行重采 $\mathbf y$）；Bos--Renes--Sprenkels 的压缩多项式 / 换 NTT 口径不同不比；pqm4 m4fstack 不掩码 5,080 B
\end{columns}
\end{frame}
\begin{frame}{masking 的下一步：换场景 \warn{笔记里查到的是 FIDO2；要不要讲你定}}
\begin{itemize}
  \item 新场景（仓库笔记，论文还没写）：\textbf{FIDO2 / WebAuthn 安全密钥}，GetAssertion 用 ML-DSA-44 签名；平台 nRF52840（64 MHz Cortex-M4F，256 KB RAM，OpenSK 用的芯片）
  \item 卖点：FIDO Server Requirements v2.3 草案把 ML-DSA-44/65/87 列为 Recommended；RFC 9964 注册了 COSE 算法号 $-48/-49/-50$
  \item 真板（只测时间和内存，不测泄漏）：n3 签名 RAM 31,944 B（仿真 31,828，+0.4\%）；起初平均 53.7 s
  \item 延迟优化：掩码 RNG 换 ChaCha20 $\to$ 24.1 s；再 bitslice A2B / 范数检查 / Decompose / $\mathbf y$ 的 B2A $\to$ \textbf{均值 6.23 s，p99 22.6 s}；目标 $\le$5 s / p99 $\le$20 s（Firefox 默认 30 s 超时）
  \item 本质不变：\textbf{masking + 省内存}，只是把「为什么要 32 KB 以内」落到一个真实产品约束上
  \item 和蚂蚁的连接：笔记里备选过「掩码 ML-KEM 解封装做认证扩展」—— 正好接他们想做的 ML-KEM
\end{itemize}
\end{frame}
\begin{frame}
\vfill\centering{\Huge\color{accent} Part 4 · HERO}\\[6pt]{\normalsize 一面已讲，备用}\vfill
\end{frame}
\begin{frame}{HERO（一面讲过，备用）}
\begin{itemize}
  \item 问题：加密只藏内容，不藏「访问了哪条记录」$\Rightarrow$ ORAM
  \item 设计：\textbf{非对称分布式 ORAM} —— 一台主服务器存加密数据库 + 一个\textbf{不共谋、资源很少的不可信 helper} 帮忙检索和 eviction；不用 FHE、不要多台全量服务器、不要 TEE
  \item 安全：主证明 semi-honest，可扩展到 malicious-with-abort
  \item 实现：OpenSSL，ElGamal 3072-bit MODP + AES-256-CTR；helper 跑在 Raspberry Pi 5（3.15 Gbps 链路）
  \item 数字：客户端带宽比 Path ORAM 最多降 \textbf{25.7$\times$}；每次访问 \textbf{0.15 s} vs Onion Ring ORAM 11.91 s；helper 峰值内存 $\approx$2 MB、代码 $\approx$105 行；新 eviction 峰值内存比前人小 2--4.9$\times$；发现 triplet eviction 参数设置有问题
  \item 简历条目里的工程细节：早期 trusted-helper 变体移植 SGX（NTL 链接进 enclave $\to$ 改 OpenSSL 路径）；ElGamal 模块化给 FPGA SmartNIC 合作者；Win/Linux/Pi 部署脚本
\end{itemize}
\src{H/intro.tex:300--334；状态：NDSS 2027 在审}
\end{frame}
\begin{frame}
\vfill\centering{\Huge\color{accent} Part 5 · threshold ML-KEM}\\[6pt]{\normalsize 组里想探索的方向}\vfill
\end{frame}
\begin{frame}{strict ML-KEM：如果让我做}
\warn{场景是我猜的：threshold 还是 masked？面试时先问一句「组里说的 ML-KEM 是多方持钥，还是抗侧信道实现？」}\\[2pt]
\begin{columns}[T]
\column{.5\textwidth}
\textbf{FIPS 203 Decaps 骨架}（难点都在这）
\begin{enumerate}
  \item $\mathbf w=v-\mathrm{NTT}^{-1}(\hat{\mathbf s}^\top\circ\mathrm{NTT}(\mathbf u))$ \hfill {\color{ok}线性 $\Rightarrow$ 本地}
  \item $m'=\mathrm{Compress}_1(\mathbf w)$ \hfill {\color{warn}舍入 = 比较 $\Rightarrow$ A2B}
  \item $(K',r')=G(m'\Vert h)$ \hfill {\color{warn}SHA3-512 on secret}
  \item $c'=\mathrm{Encrypt}(ek,m',r')$：CBD 采样（SHAKE $\to$ B2A），$A\mathbf r$ 线性，$\mathrm{Compress}_{d_u,d_v}$ 又是舍入
  \item $\bar K=J(z\Vert c)$；\textbf{$c\stackrel{?}{=}c'$}，不能泄露走了哪条分支（implicit rejection）
\end{enumerate}
\textbf{$q=3329$}（12 bit），比 ML-DSA 小，域更小 $\Rightarrow$ MAC soundness 问题更突出
\column{.5\textwidth}
\textbf{threshold ML-KEM}：SHIELD 打法直接迁移 —— SPDZ 做 1、4 的线性部分；WRK17 做 Keccak、Compress、比较；daBits 桥接；结果用 oblivious MUX 选 $K'$/$\bar K$\\[3pt]
\textbf{masked ML-KEM}：同样的积木 —— masked Compress（A2B）、masked Keccak、masked CBD（B2A）、\textbf{masked comparison}（历史上最容易出错的一块）、FO 的 plaintext-checking-oracle 侧信道\\
\src{凭记忆的前作：Bos et al. ``Masking Kyber'' TCHES'21；Heinz et al. first-order masked Kyber on M4 (ePrint 2022/058)；Bhasin et al. masked comparison 攻防 TCHES'21 —— \warn{引用前核对}}\\[3pt]
\textbf{我能带过去的经验}：bit-exact 合规的工程纪律；A2B$_q$（Compress 就是 HighBits 的翻版，$q$ 只有 12 bit 更便宜）、B2A（CBD 采样）；4 份 TI Keccak 直接复用；「只合并 1 个判定位」用在 FO 比较；R1/R2 barrier 与编译器无关可直接迁移\\
\textbf{威胁模型差别}：KEM 面对\textbf{选择密文}，FO 判定位和比较 gadget 是主攻击面；ML-DSA 是已知消息
\end{columns}
\end{frame}
\begin{frame}
\vfill\centering{\Huge\color{accent} Part 6 · 概念讲解}\\[6pt]{\normalsize 我没讲清 / 没意识到的，逐个弄明白}\vfill
\end{frame}
\begin{frame}{E1 · ISW 是什么（Ishai--Sahai--Wagner, ``Private Circuits: Securing Hardware against Probing Attacks'', CRYPTO 2003）}
\begin{columns}[T]
\column{.5\textwidth}
\textbf{要解决的问题}：攻击者能在芯片上「探」到 $t$ 根导线的值（\textbf{$t$-probing model}，侧信道的理论抽象）。怎么让他探 $t$ 根线也学不到秘密？\\[3pt]
\textbf{做法}：把每个比特拆成 $d+1$ 份，$x=x_0\oplus x_1\oplus\dots\oplus x_d$；任意 $d$ 份都是均匀随机 $\Rightarrow$ 探 $d$ 根线什么都看不到\\[3pt]
\begin{itemize}
  \item \textbf{XOR / 线性运算}：逐份算，$c_i=a_i\oplus b_i$，免费
  \item \textbf{AND 难}：$a\wedge b=\bigoplus_{i,j}a_ib_j$ 有交叉项 $a_ib_j$，单独算出来就把两份放一起了 $\Rightarrow$ 需要随机数把交叉项遮住 = \textbf{ISW AND gadget}
\end{itemize}
\textbf{术语}：gadget = 一个带掩码的小电路模块；refresh = 用新随机数重新拆份
\column{.5\textwidth}
\textbf{ISW AND（$d+1$ 份）}
\begin{enumerate}
  \item 对每对 $i<j$ 采随机比特 $r_{ij}$
  \item $r_{ji}=(r_{ij}\oplus a_ib_j)\oplus a_jb_i$ \quad（\alert{括号顺序不能换}）
  \item $c_i=a_ib_i\oplus\bigoplus_{j\neq i}r_{ij}$
\end{enumerate}
随机数 $d(d+1)/2$ 个；乘法 $(d+1)^2$ 次\\[3pt]
\textbf{2 份例子}：$c_0=a_0b_0\oplus r$，\quad $c_1=a_1b_1\oplus\big((r\oplus a_0b_1)\oplus a_1b_0\big)$\\[3pt]
\textbf{和 MPC 的联系（面试金句）}：「这就是 GMW 协议里的 AND：每个份额看成一个参与方，交叉项 $a_ib_j$ 由 $i$、$j$ 两方用共享随机数 $r_{ij}$ 遮住后再合并。」\\[3pt]
\src{原论文用 $2t+1$ 份证明；Rivain--Prouff (CHES 2010) 证明 $d+1$ 份的 ISW AND 是 $d$-probing 安全（并修了 refresh）\warn{出处凭记忆}}
\end{columns}
\end{frame}
\begin{frame}{E2 · masking、MPC、MPC-in-the-head 三者关系}
{\scriptsize\renewcommand{\arraystretch}{1.2}
\begin{tabularx}{\textwidth}{@{}p{2.2cm}YYY@{}}\toprule
 & \textbf{MPC}（如 SPDZ、WRK17）& \textbf{masking}（如 ISW、TI）& \textbf{MPC-in-the-head}（IKOS07）\\\midrule
目的 & 多方联合计算，谁都不看到别人输入 & 一块芯片抗侧信道 & 构造 \textbf{zero-knowledge proof} \\
「参与方」 & 真实的不同机器 & 同一芯片上的 $d+1$ 个份额 & prover 脑中模拟的 $n$ 个虚拟方 \\
敌手 & 腐化至多 $t$ 方（半诚实 / 恶意）& 探测至多 $d$ 个中间值（probing / noisy leakage）& verifier 要求打开 $t$ 个虚拟方的 view 来检查 \\
能「打开」值吗 & 能：打开被随机数遮住的值（Beaver 的 $d,e$；A2B 打开 $x-r$）& 基本不能：芯片上任何组合值都可能被探到 & 打开 $t$ 个 view \\
完整性 & 恶意安全靠 IT-MAC & 只管隐私，抗故障注入要另做 & 靠 commit + 随机打开 \\
例子 & SHIELD、FLOCK & 我的 masked ML-DSA & Picnic、Ligero / Longfellow、ZKBoo \\\bottomrule
\end{tabularx}}
\vspace{.4em}

\textbf{所以「masking 是 local MPC 吗？」} —— \textbf{思想上是}（ISW AND $\approx$ GMW AND），但模型更弱的地方正是我论文的发现：
\begin{itemize}
  \item MPC 的各方物理隔离；芯片上的「各方」\alert{共用寄存器、总线、操作数端口}，前后两个值的 Hamming distance 会漏（transition），硬件里还有 glitch
  \item 所以「每一步只碰一份」的 ISW 证明在\textbf{值模型}成立，到了编译后的二进制上会被总线打破 $\Rightarrow$ R1/R2、bus-and-port barrier
\end{itemize}
\src{IKOS07 = Ishai, Kushilevitz, Ostrovsky, Sahai, ``Zero-Knowledge from Secure Multiparty Computation'', STOC 2007}
\end{frame}
\begin{frame}{E3 · masking 圈的术语（一面那位提到 ISW，老板可能接着问）}
\begin{columns}[T]
\column{.5\textwidth}
\begin{itemize}
  \item \textbf{$d$-th order masking}：$d+1$ 份，抗 $d$ 个探针
  \item \textbf{value model vs transition model}：只看单个值 vs 也看「前一个值 $\oplus$ 后一个值」（寄存器 / 总线被覆写时的 HD 泄漏）。Balasch et al.：$n$ 份在 transition 下只剩 $\lfloor(n-1)/2\rfloor$ 阶 $\Rightarrow$ 我们检查路径用 3 份
  \item \textbf{glitch}：硬件组合逻辑里信号短暂的毛刺，会让一根线「看到」多份 $\Rightarrow$ 硬件 masking 要 TI / DOM
  \item \textbf{NI / SNI}（Non-Interference / \textbf{Strong} Non-Interference, Barthe et al. EUROCRYPT 2016）：gadget 的探针能用多少个输入份模拟。SNI = 只与\textbf{内部}探针数有关，与输出探针无关 $\Rightarrow$ 多个 gadget 可以安全\textbf{组合}
  \item \textbf{PINI}（Probe-Isolating Non-Interference, Cassiers--Standaert TCHES 2020）：更容易组合的变体
\end{itemize}
\column{.5\textwidth}
\begin{itemize}
  \item \textbf{TI}（Threshold Implementation, Nikova--Rechberger--Rijmen 2006）：每个输出份只用到部分输入份（\textbf{non-completeness}）$\Rightarrow$ 即使有 glitch 也不会一次碰到全部份额。「threshold」\alert{不是}门限密码
  \item \textbf{DOM}（Domain-Oriented Masking, Groß--Mangard--Korak 2016）：每个份额一个「域」，跨域乘积加随机数后再寄存
  \item \textbf{changing of the guards}（Daemen, CHES 2017）：Keccak $\chi$ 的 TI 不用每轮新随机数的技巧
  \item \textbf{refresh}：用新随机数重拆份；我们用的 RP10 式线性 refresh 一阶均匀，但\textbf{不是 SNI}
  \item \textbf{TVLA}（Test Vector Leakage Assessment）：固定输入 vs 随机输入两组 trace 做 Welch $t$-test，$|t|$ 超阈值 = 检出泄漏
  \item \textbf{CPA}（Correlation Power Analysis）/ \textbf{template attack}：真正去恢复密钥的攻击
\end{itemize}
\end{columns}
\end{frame}
\begin{frame}{E4 · SPDZ 怎么做到 malicious，以及 MPC 里的 A2B}
\begin{columns}[T]
\column{.5\textwidth}
\textbf{IT-MAC}（information-theoretic MAC）
\begin{itemize}
  \item 全局 MAC key $\alpha$ 被加法拆给各方，谁也不知道
  \item 每个秘密 $x$：各方持 $(x_i,m_i)$，$\sum x_i=x$，$\sum m_i=\alpha x$
  \item 恶意方把 $x$ 改成 $x+\delta$，还得把 MAC 改成 $\alpha(x+\delta)$ —— 不知道 $\alpha$，只能猜中概率 $1/p$ $\Rightarrow$ \textbf{MAC check} 时被抓 $\Rightarrow$ abort
  \item 所以 $p$ 小（$q\approx2^{23}$）时要多个 $\alpha$ 或扩域
\end{itemize}
\textbf{Beaver triple}：预处理好 $([a],[b],[c=ab])$；算 $[xy]$ 时打开 $d=x-a,e=y-b$（均匀随机，不泄露），$[xy]=[c]+d[b]+e[a]+de$ $\Rightarrow$ online 一轮
\column{.5\textwidth}
\textbf{MPC 里的 A2B}（对照 masking 的 A2B）
\begin{itemize}
  \item 预处理一个 \textbf{daBit / edaBits}：同一个随机数 $r$ 同时有算术份 $[r]_A$ 和逐比特布尔份 $[r_0],\dots,[r_{k-1}]_B$
  \item 打开 $c=x-r$（被 $r$ 遮住，公开无害）
  \item 布尔电路里算 $x=c+r$（加法器 + mod $p$ 修正）
  \item edaBits = Escudero et al. CRYPTO 2020，「一次拿一串比特」比逐比特 daBit 省
\end{itemize}
\textbf{对照}：MPC 能「打开 $x-r$」，所以 A2B 只要一个公开数 + 一个加法器；芯片上不能打开，只能用 SecAdd 在份额上整体算 $\Rightarrow$ masking 的 A2B 更贵\\[3pt]
\textbf{abort 的种类}：security-with-abort（我们）< identifiable abort（能指出谁作弊）< robustness / guaranteed output delivery
\end{columns}
\end{frame}
\begin{frame}{E5 · 为什么拒绝采样让 threshold 化这么难（malicious 的根本理由）}
\begin{columns}[T]
\column{.5\textwidth}
\textbf{Fiat--Shamir with aborts}（ML-DSA 的核心）
\begin{itemize}
  \item $\mathbf z=\mathbf y+c\,\mathbf s_1$；$\mathbf y$ 是随机 mask
  \item 若 $\|\mathbf z\|_\infty\ge\gamma_1-\beta$ 就\textbf{拒绝重来}
  \item 这样\textbf{被接受}的 $\mathbf z$ 分布与 $\mathbf s_1$ 无关；\textbf{被拒绝}的 $\mathbf z$ 是「往 $c\mathbf s_1$ 方向偏」的样本 $\Rightarrow$ 泄露 $\mathbf s_1$
  \item 期望重试：44/65/87 = 4.25 / 5.1 / 3.85 次
\end{itemize}
\textbf{threshold 的难点}：多方要在\textbf{不暴露候选 $\mathbf z$} 的前提下共同决定拒不拒绝。各路线的回答：
\begin{itemize}
  \item SHIELD：整个检查在 MPC 里，只公开 1 个 reject bit
  \item Celi：短份额 + 各方本地拒绝（改分布）
  \item Quorus：加噪声 $e_w$，使被拒的 $w_1$ 可安全公开
\end{itemize}
\column{.5\textwidth}
\textbf{为什么半诚实不够（3 个具体攻击）}
\begin{enumerate}
  \item \textbf{篡改判定}：恶意方在 $\mathbf z$ 的份或 reject bit 上加 $\delta$ $\Rightarrow$ 本该拒绝的 $\mathbf z$ 被打开 $\Rightarrow$ 泄露 $\mathbf s_1$
  \item \textbf{selective failure}：让「是否 abort」依赖某个秘密比特；因为拒绝本来就常见，人为 abort 和正常拒绝分不出 $\Rightarrow$ 多次会话逐比特提取。$\Rightarrow$ 影响后续行为的值必须\textbf{先认证再用}
  \item \textbf{nonce 偏置 / rushing}：看到别人的贡献再选自己的 $\Rightarrow$ $\mathbf y$ 有偏（Celi 用「先 commit 哈希再公开」防）
\end{enumerate}
\textbf{一句话}：「threshold 签名里持钥方本身就是敌手，而 ML-DSA 的安全恰恰靠『被拒的东西永远不出现』—— 半诚实模型保证不了这一点。」
\end{columns}
\end{frame}
\begin{frame}{E6 · 几种 secret sharing 一页看懂}
{\scriptsize\renewcommand{\arraystretch}{1.25}
\begin{tabularx}{\textwidth}{@{}p{2.6cm}YYp{3.2cm}@{}}\toprule
 & 怎么分 & 优点 & 代价 / 谁用 \\\midrule
\textbf{additive}（$n$-of-$n$）& $x=\sum x_i$ & 最简单；配 IT-MAC 就是 SPDZ，容忍 $n-1$ 恶意 & 少一方就签不了；SHIELD \\
\textbf{Shamir}（$t$-of-$n$）& $x=f(0)$，$f$ 是 $t-1$ 次多项式，各方拿 $f(i)$ & 任意 $t$ 方能重构；可扩展到很多方 & 份额在 $\mathbb Z_q$ 均匀、Lagrange 系数大 $\Rightarrow$ 不能本地做拒绝；乘法要 honest majority；Quorus \\
\textbf{replicated (RSS)} & 拆成多块加法份，每块给一组人 & 重构系数 0/1；乘法本地化 & 份数 $\binom{N}{T-1}$ 随 $N$ 指数涨 \\
\textbf{short RSS} & RSS，但每块都是\textbf{短}的 ML-DSA 秘密 & 每方的部分秘密本身就短 $\Rightarrow$ 能本地拒绝采样 & $N\le6$；Celi / Mithril \\\bottomrule
\end{tabularx}}
\vspace{.5em}

\textbf{一句话串起来}：「additive + MAC 给最强的威胁模型但只能 $n$-of-$n$；Shamir 给 $t$-of-$n$ 但要 honest majority；short RSS 让拒绝采样可以本地做，但份数爆炸。我们选第一种是因为 FPKI 要 $n$-of-$n$ 且容忍 $n-1$。」
\end{frame}
\begin{frame}{术语别说错 + 数字速查}
\begin{columns}[T]
\column{.5\textwidth}
\textbf{术语}
\begin{itemize}
  \item \alert{masking $\neq$ MPC-in-the-head}。MPC-in-the-head 是 IKOS07 用「脑中模拟 MPC」构造 ZK 证明（Picnic、Longfellow 那类）。
        masking 应该说：\textbf{「单设备上的 MPC / ISW private circuits，probing model 下每个中间值拆成 $d+1$ 份」}
  \item A2B / B2A：arithmetic $\leftrightarrow$ Boolean masking 转换
  \item short RSS（Celi）vs Shamir（Quorus）vs additive + IT-MAC（SPDZ）
  \item security-with-abort / identifiable abort / robustness
  \item functional interchangeability vs algorithmic compliance
\end{itemize}
\column{.5\textwidth}
\textbf{数字}
\begin{itemize}
  \item ML-DSA $q=8380417=2^{23}-2^{13}+1$；ML-KEM $q=3329$
  \item 期望重试：44/65/87 = 4.25 / 5.1 / 3.85
  \item SHIELD：2--5 方；单次迭代 198--536 ms；23 轮；7.7--14.7 MB/CA
  \item WRK17 online 4 轮；SampleInBall 890K$\to$355K（相对最初编码）
  \item Celi：$\le$6 方，3 轮/次，LAN < 20 ms
  \item Quorus：honest maj.，测到 63 方
  \item masking：峰值 RAM $-37.1/-39.5/-41.9\%$；44: 50,612 $\to$ 31,828 B
  \item HERO：带宽 25.7$\times$，0.15 s/访问
  \item FLOCK：9 轮，4 方 25 ms
\end{itemize}
\end{columns}
\end{frame}

\end{document}
